\documentclass{article}

\usepackage{a4wide}
\usepackage{amsmath,amssymb,amsthm}
\usepackage{natbib}
\usepackage{booktabs}
\usepackage{tikz}

\newtheorem{example}{Example}
\newtheorem{theorem}{Theorem}
\newtheorem{definition}{Definition}
\newtheorem{observation}{Observation}
\newtheorem{proposition}{Proposition}
\newtheorem{conjecture}{Conjecture}
\newtheorem{corollary}{Corollary}
\newtheorem{lemma}{Lemma}

%Shortcuts
\newcommand{\match}{\mu}
\newcommand{\linearorders}{\mathcal{L}}

%TikZ commands
\newcommand*\encircled[1]{\tikz[baseline=(char.base)]{
            \node[shape=circle,draw,inner sep=0.25pt,minimum size=1em] (char) {\vphantom{Ag}#1};}}
\newcommand*\framed[1]{\tikz[baseline=(char.base)]{
            \node[draw,inner sep=0.25pt,minimum size=1em] (char) {\vphantom{Ag}#1};}}
            

\title{Possible and Necessary Popularity in House Allocation}

\begin{document}

\maketitle

\section{Preliminaries}

We are given a set of $n$ agents $N = \{1,\dots,n\}$, and a set of $n$ items $O=\{o_1,\dots,o_n\}$.
We assume that each agent $i \in N$ has strict ordinal preferences on the items, represented by a linear order $\succ_i$ over $O$.
The preference profile $\succ$ is the collection of all agents' preferences, i.e., $\succ=(\succ_1,\dots,\succ_n)$.
The set of all possible linear orders over $O$ is denoted by $\linearorders(O)$.
A matching $\match:N\to O$ is a bijection assigning to every agent $i\in N$ an item in $O$ denoted by $\match(i)$.

A matching $\match$ is \emph{Pareto-optimal} if there is no other matching that is at least as good for all agents and strictly better for at least one agent.
We are interested by a refinement of Pareto-optimality called \emph{popularity}: a matching is popular if there is no other matching that is more popular, i.e., that is preferred by more agents.

\begin{definition}[Popularity]
A matching $\match$ is popular if $|\{i\in N: \match(i)\succ_i \match'(i)\}|\geq |\{i\in N: \match'(i)\succ_i \match(i)\}|$, for every matching $\match'$.
\end{definition}

A popular matching is trivially Pareto-optimal but is not guaranteed to exist.
The following characterization of popularity due to \citet{abraham2007popular} allows to decide in polynomial time whether a popular matching exists.

\begin{lemma}[\citet{abraham2007popular}]\label{lem:charact-pop}
A matching is popular iff every agent is assigned either her most preferred item or her most preferred item among those that are not ranked first by any agent.
\end{lemma}

We suppose that the agents may communicate partial preferences, i.e., each agent $i\in N$ expresses a strict partial order $P_i$ over $O$.
The partial preference profile is denoted by $P=(P_1,\dots,P_n)$.
An instance of house allocation with incomplete preferences is thus given by $\langle N,P \rangle$.

To reconstruct the complete preferences, we will consider possible completions of the partial preference profile.
A linear order $\succ_i$ is a linear extension of $P_i$ if it is compatible with $P_i$, i.e., for all items $o,o'\in O$ if $o P_i o'$ then $o \succ_i o'$. 
A preference profile $\succ$ is a possible completion of a partial preference profile $P$ if $\succ_i$ is a linear extension of $P_i$, for every $i\in N$.
%The set of all possible completions of a partial preference $P_i$ is denoted by $Lin(P_i)$.

The definition of popularity can then be adapted to partial preferences by considering possible or necessary outcomes with respect to the possible preference completions.

\begin{definition}[Possible popularity]
A matching is possibly popular if it is popular in at least one completion of the partial preference profile.
\end{definition}

\begin{definition}[Necessary popularity]
A matching is necessarily popular if it is popular in all completions of the partial preference profile. 
\end{definition}

Let us illustrate these two notions in the following example.

\begin{example}
Let us consider an instance of house allocation with incomplete preferences with four agents who have the following partial preferences:
\begin{center}
\begin{tabular}{ll}
$1:$ & $a\succ b$ \\
$2:$ & $b\succ c$ \\
$3:$ & $a\succ c$ \\
$4:$ & $a\succ d\succ b$
\end{tabular}
\end{center}
There exists the following completion of the partial preference profile where there are exactly two popular matchings: the framed and the encircled matchings.
\begin{center}
\begin{tabular}{*{8}{c}}
$1:$ & $\framed{\encircled{a}}$ & $\succ$ & $b$ & $\succ$ & $c$ & $\succ$ & $d$ \\
$2:$ & $\framed{b}$ & $\succ$ & $a$ & $\succ$ & $\encircled{c}$ & $\succ$ & $d$ \\
$3:$ & $\encircled{b}$ & $\succ$ & $a$ & $\succ$ & $\framed{c}$ & $\succ$ & $d$ \\
$4:$ & $a$ & $\succ$ & $\framed{\encircled{d}}$ & $\succ$ & $c$ & $\succ$ & $b$
\end{tabular}
\end{center}
It follows that these two matchings are possibly popular in the initial partial preference profile. 
However, none of them is necessarily popular as they would not be popular in the completion where the partial preferences of agent 1 would be extended in such a way: $d \succ_1 a \succ_1 b$.  
Since these matchings are the only popular matchings in the previously presented preference completion, in general, there are no necessarily popular matchings in the initial partial preference profile.
\end{example}

We can then investigate the following computational problems:
\smallbreak
\begin{tabular}{ll}
\toprule
\multicolumn{2}{c}{\textsc{ExistPossibPop}/\textsc{ExistNecessPop}} \\
\midrule
Instance: & House allocation instance with incomplete preferences $(N,P)$ \\
Question: & Does there exist a possibly/necessarily popular matching? \\
\bottomrule
\end{tabular}

\smallbreak
\begin{tabular}{ll}
\toprule
\multicolumn{2}{c}{\textsc{VerifPossibPop}/\textsc{VerifNecessPop}} \\
\midrule
Instance: & House allocation instance with incomplete preferences $(N,P)$, matching $\match$ \\
Question: & Is $\match$ possibly/necessarily popular? \\
\bottomrule
\end{tabular}

\bibliographystyle{plainnat}
\bibliography{biblio}

\end{document}